%% ma: L12-B06-0004
%% lop: 12
%% chuong: 2
%% bai: 6
%% dang: 12.6.3
%% loai: TN4
%% muc: TH
%% dap_an: "D"
%% nguon: "(NBV)-ĐỀ SỐ 4-MỖI NGÀY 1 ĐỀ THI 2025.pdf (D04, MỖI NGÀY 1 ĐỀ - Đề số 4), Phần I Câu 3"
%% kien_thuc: "Bài 6 (tổng các vectơ, trung điểm, trọng tâm; quy tắc trọng tâm)"
%% trang_thai: chinh-thuc
%% ghi_chu: "Cùng dạng 12.6.3 với L12-B06-0003 (đây là câu chọn khẳng định đúng, dùng thêm trọng tâm). Hình gốc chỉ minh họa nên không vẽ lại. Đáp án gốc D đúng (đã kiểm bằng số ngẫu nhiên)."
\begin{ex}
Cho tứ diện $ABCD$. Gọi $I,J$ theo thứ tự là trung điểm của các đoạn thẳng $AB,CD$ và $G$ là trọng tâm của tam giác $BCD$. Khẳng định nào sau đây đúng?
\choice{$3\overrightarrow{BG}+2\overrightarrow{BJ}=\vec0$}{$\overrightarrow{IB}+\overrightarrow{IC}+\overrightarrow{ID}=3\overrightarrow{IJ}$}{$\overrightarrow{GA}+\overrightarrow{GB}+\overrightarrow{GC}=\vec0$}{\True $\overrightarrow{AB}+\overrightarrow{AC}+\overrightarrow{AD}=3\overrightarrow{AG}$}
\loigiai{$G$ là trọng tâm tam giác $BCD$ nên $\overrightarrow{GB}+\overrightarrow{GC}+\overrightarrow{GD}=\vec0$, suy ra $\overrightarrow{AB}+\overrightarrow{AC}+\overrightarrow{AD}=3\overrightarrow{AG}$ (phương án D đúng).
Các phương án còn lại sai: $\overrightarrow{BG}=\frac23\overrightarrow{BJ}$ nên $3\overrightarrow{BG}-2\overrightarrow{BJ}=\vec0$ (A sai); $\overrightarrow{IB}+\overrightarrow{IC}+\overrightarrow{ID}=3\overrightarrow{IG}\neq3\overrightarrow{IJ}$ (B sai); $\overrightarrow{GA}+\overrightarrow{GB}+\overrightarrow{GC}=\overrightarrow{GA}-\overrightarrow{GD}=\overrightarrow{DA}\neq\vec0$ (C sai).}
\end{ex}
